\documentclass[12pt, a4paper]{report}

% Packages
\usepackage[utf8]{inputenc}
\usepackage[T1]{fontenc}
\usepackage[english]{babel}
\usepackage{amsmath, amssymb, amsthm, amsfonts}
\usepackage{geometry}
\usepackage{hyperref}
\usepackage{bookmark}
\usepackage{enumitem}
\usepackage{fancyhdr}
\usepackage{titlesec}
\usepackage{mathrsfs}
\usepackage{xcolor}
\usepackage{tikz}
\usepackage{pgfplots}
\pgfplotsset{compat=1.18}

% Geometry setup
\geometry{
    top=2.5cm,
    bottom=2.5cm,
    left=2.5cm,
    right=2.5cm
}
\setlength{\headheight}{30pt}

% Hyperref setup
\hypersetup{
    colorlinks=true,
    linkcolor=blue,
    filecolor=magenta,      
    urlcolor=cyan,
    pdftitle={Stochastic Finance and Risk Modelling - Chapter 6},
    pdfauthor={Laurence Carassus and Gaoyue Guo},
}

% Theorem environments
\newtheorem{theorem}{Theorem}[chapter]
\newtheorem{proposition}[theorem]{Proposition}
\newtheorem{lemma}[theorem]{Lemma}
\newtheorem{corollary}[theorem]{Corollary}

\theoremstyle{definition}
\newtheorem{definition}{Definition}[chapter]
\newtheorem{example}{Example}[chapter]
\newtheorem{remark}[theorem]{Remark}

% Layout/Style
\pagestyle{fancy}
\fancyhf{}
\rhead{\leftmark}
\lhead{}
\cfoot{\thepage}

\title{\textbf{Stochastic Finance and Risk Modelling} \\ \large Chapter 6: The Binomial Model}
\author{Based on slides by Ioane Muni-Toke \\ Laurence Carassus and Gaoyue Guo \\ \textit{Expanded Study Resource}}
\date{\today}

\begin{document}

\maketitle

\tableofcontents

\chapter*{Introduction}
\addcontentsline{toc}{chapter}{Introduction}
This document serves as a comprehensive study guide for the course "Stochastic Finance and Risk Modelling", functioning as a complete substitute for the lecture slides related to Chapter 6: "The Binomial Model".

The celebrated Cox-Ross-Rubinstein (CRR) binomial model, introduced in 1979, is one of the foundational frameworks in modern financial mathematics. Despite its simplicity, it illustrates the core mechanics of arbitrage pricing theory in discrete time. In this chapter, we formally define the binomial market setup for a risky and a riskless asset. We then formulate the probability space and natural filtration associated with the price process, verify the conditions for no-arbitrage, and compute the unique equivalent martingale measure. Finally, we explore derivative pricing through backward induction, demonstrate how any European derivative can be replicated, and establish why the binomial model is intrinsically a complete market. 

To aid intuition, rigorous mathematical derivations, detailed explanations, and worked-out proofs of the main theorems have been fully developed to bridge conceptual gaps left by the slides.

\chapter{The Binomial Model}

\section{Description and Market Setup}

The binomial model, largely popularized by the seminal 1979 paper by John Cox, Stephen Ross, and Mark Rubinstein, represents one of the most elegant and fundamental discrete-time models in financial mathematics. Its power lies in its simplicity: at any given time step, the asset price can only move to one of two possible states (up or down). While initially designed as a discretized approximation of the continuous-time Black-Scholes-Merton model, the binomial model stands on its own as a robust logical framework for illustrating the core principles of no-arbitrage, risk-neutral valuation, and dynamic hedging.

We consider a discrete-time financial market running up to a deterministic finite horizon $T \in \mathbb{N}$. That is, trading occurs at times $t \in \{0, 1, \dots, T\}$. For the sake of simplicity and illustrative clarity, we restrict our market to contain exactly two assets: one riskless asset (e.g., a bank account or a zero-coupon bond) and $d=1$ risky asset (e.g., a stock).

\subsection{The Riskless Asset}
The riskless asset, denoted by $S^0 = (S^0_t)_{t=0,1,\dots,T}$, grows at a constant deterministic interest rate $r \geq 0$. We normalize its initial value to $S^0_0 = 1$. The value of the riskless asset at time $t$ is therefore given by the compounded growth:
\begin{equation}
    S^0_t = (1+r)^t, \quad \text{for } t = 0, 1, \dots, T.
\end{equation}
Consequently, the one-period simple return on the riskless asset at any time $t \in \{1, \dots, T\}$ is deterministic and equal to $r$:
\begin{equation}
    R^0_t := \frac{S^0_t - S^0_{t-1}}{S^0_{t-1}} = \frac{(1+r)^t - (1+r)^{t-1}}{(1+r)^{t-1}} = \frac{(1+r)^{t-1}((1+r) - 1)}{(1+r)^{t-1}} = r.
\end{equation}

\subsection{The Risky Asset}
The risky asset is denoted by a stochastic process $S = (S_t)_{t=0,1,\dots,T}$. Its initial price is a known constant $S_0 = s > 0$. Its value at any subsequent time step is uncertain and modeled mathematically as a random variable. The one-period simple return of the risky asset at time $t$ is defined as:
\begin{equation}
    R_t := \frac{S_t - S_{t-1}}{S_{t-1}}.
\end{equation}
Equivalently, the price transitions recursively according to:
\begin{equation}
    S_t = S_{t-1}(1 + R_t).
\end{equation}
The defining characteristic of the binomial model is the assumption that the one-period returns $R_t$ can only take one of two possible values at any step. Specifically, we assume that with probability 1:
\begin{equation}
    R_t \in \{d, u\}, \quad \text{for } t=1, \dots, T
\end{equation}
where $d$ and $u$ are real numbers representing the "down" and "up" scenarios, respectively. To ensure that the asset price remains strictly positive ($S_t > 0$ almost surely) and that the "up" movement is strictly greater than the "down" movement, we impose the strict inequality:
\begin{equation}
    -1 < d < u.
\end{equation}

\section{Probability Space and Information Structure}

To formally ground the binomial model in measure-theoretic probability, we must define the triplet $(\Omega, \mathcal{F}, \mathbb{P})$ and the corresponding filtration $\mathbb{F} = (\mathcal{F}_t)_{t=0}^T$.

\subsection{The Sample Space}
Since there are $T$ trading periods and exactly 2 possible outcomes for the return in each period, the sample space $\Omega$ simply consists of all possible sequences of length $T$ composed of the letters $u$ and $d$. Thus,
\begin{equation}
    \Omega := \{d, u\}^T.
\end{equation}
A typical element $\omega \in \Omega$ represents a complete path of the market up to maturity $T$, denoted as $\omega = (\omega_1, \omega_2, \dots, \omega_T)$. Since $\Omega$ is finite, it contains $2^T$ distinct paths. The natural choice for the $\sigma$-algebra $\mathcal{F}$ is the power set, which contains all possible subsets of $\Omega$:
\begin{equation}
    \mathcal{F} := 2^{\Omega}.
\end{equation}

\subsection{Random Variables and Filtration}
We define the single-period return at time $t$ by evaluating the coordinate projection mapping $R_t : \Omega \to \{d, u\}$ given by:
\begin{equation}
    R_t(\omega) := \omega_t, \quad \text{for all } \omega = (\omega_1, \dots, \omega_T) \in \Omega.
\end{equation}
The strictly positive risky asset price $S_t$ is therefore explicitly defined as a cumulative product of these returns:
\begin{equation}
    S_t(\omega) = S_0 \prod_{k=1}^t (1 + R_k(\omega)) = S_0 \prod_{k=1}^t (1 + \omega_k).
\end{equation}

Information is revealed dynamically over time. At time $t=0$, nothing has happened yet, so the initial $\sigma$-algebra models total uncertainty to the observer outside of deterministic initial values: $\mathcal{F}_0 = \{\emptyset, \Omega\}$. For any $t \in \{1, \dots, T\}$, the information available incorporates all the historical returns observed up to time $t$. The filtration is defined to represent this information accumulation:
\begin{equation}
    \mathcal{F}_t := \sigma(R_1, \dots, R_t).
\end{equation}

\begin{remark}
    By construction, $\mathcal{F}_t = \sigma(S_1, \dots, S_t)$. To verify this, notice that $S_k = S_{k-1}(1+R_k)$. If we observe the returns $\{R_1, \dots, R_t\}$, we can compute the path of prices $\{S_1, \dots, S_t\}$. Conversely, if we strictly observe prices $\{S_1, \dots, S_t\}$ along with the deterministic initial price $S_0$, we can deduce the exact historical returns recursively using $R_k = \frac{S_k}{S_{k-1}} - 1$. Because the transformation is bijective, they generate functionally identical $\sigma$-algebras. Furthermore, $\mathcal{F}_T = \mathcal{F}$, confirming that all systemic uncertainty is strictly resolved at maturity $T$.
\end{remark}

\subsection{Reference Probability Measure}
We endow the finite measurable space $(\Omega, \mathcal{F})$ with an objective probability measure $\mathbb{P} \in \mathscr{P}(\Omega, \mathcal{F})$. Our only baseline assumption on the physical measure $\mathbb{P}$ is that every path is theoretically possible. Thus, we strongly enforce $\mathbb{P}(\{\omega\}) > 0$ for every specific path $\omega \in \Omega$.

\section{No-Arbitrage and Completeness}

The principle of no-arbitrage rests as the foundational qualitative assumption in modern finance. It broadly guarantees that it is impossible to orchestrate a sophisticated trading strategy to yield a strictly positive riskless profit from zero initial capital investment.

\begin{theorem}[No-Arbitrage Condition in the Binomial Model]
    The binomial model is strictly arbitrage-free if and only if the prevailing interest rate lies strictly between the boundaries of the downward and upward asset returns:
    \begin{equation}
        d < r < u. \label{eq:no_arb_condition}
    \end{equation}
\end{theorem}

\begin{proof}
    We can derive this intuition quite simply from a one-period perspective. If $r \leq d < u$, the risky asset deterministically yields a return at least equal to, and frequently strictly greater than, the riskless asset. An aggressive investor could borrow limitless cash at rate $r$ to blindly buy the stock, generating a strictly positive riskless windfall. Specifically, taking a loan of $\$1$ at time $t-1$ requires repaying $(1+r)$ at time $t$. The purchased stock is physically worth $(1+R_t)$ at time $t$. The net profit resolves to $1+R_t - (1+r) = R_t - r \geq d - r \geq 0$, and with $R_t = u > r$, the specific profit stretches strictly positive. Since $\mathbb{P}(R_t=u) > 0$, an arbitrage violently materializes.
    
    Conversely, if $d < u \leq r$, the riskless asset rigidly dominates the risky asset. Short-selling the stock to strategically invest the proceeds into the bank account secures an identical absolute arbitrage. The only mathematical mechanism capable of eliminating these structural disparities is enforcing tight boundaries around the riskless return between the two possible states of the risky return: $d < r < u$.
\end{proof}

When the structural condition of no-arbitrage seamlessly holds, the discrete-time First Fundamental Theorem of Asset Pricing guarantees the abstract existence of an Equivalent Martingale Measure (EMM) $\mathbb{Q}$. In the deeply restricted binomial framework, this EMM exhibits a highly analytical definition. Furthermore, the abstract Second Fundamental Theorem asserts structural market completeness strictly derived from the absolute uniqueness of $\mathbb{Q}$.

\begin{theorem}[Completeness and the Risk-Neutral Measure]
    If $d < r < u$ securely holds, the binomial model represents a \textbf{complete} market. Additionally, there exists a strictly \textbf{unique} equivalent martingale measure $\mathbb{Q}$, defined functionally by the explicit property that the coordinate returns $R_1, \dots, R_T$ are aggressively independent and identically distributed (i.i.d.) under $\mathbb{Q}$ strictly satisfying:
    \begin{equation}
        \mathbb{Q}[R_t = u] = \frac{r - d}{u - d} := q \in (0, 1), \quad \text{for every } t = 1, \dots, T.
    \end{equation}
    As a consequential corollary, the underlying price process $(S_t)_{t=0, \dots, T}$ forms a structural Markov process rigidly under $\mathbb{Q}$.
\end{theorem}

\begin{proof}
    By the canonical definition, $\mathbb{Q}$ acts as an EMM if and only if the underlying discounted price process evaluated as $\widetilde{S}_t := S_t / (1+r)^t$ behaves exactly as a martingale under $\mathbb{Q}$. Translated mathematically:
    \begin{equation}
        \mathbb{E}^\mathbb{Q} \left[ \widetilde{S}_t \mid \mathcal{F}_{t-1} \right] = \widetilde{S}_{t-1}.
    \end{equation}
    Directly substituting the recursive definition into the expectation:
    \begin{equation}
        \mathbb{E}^\mathbb{Q} \left[ \frac{S_{t-1}(1+R_t)}{(1+r)^t} \mid \mathcal{F}_{t-1} \right] = \frac{S_{t-1}}{(1+r)^{t-1}}.
    \end{equation}
    Because structural information up to time $t-1$ precisely determines $S_{t-1}$ (i.e., it holds as strongly $\mathcal{F}_{t-1}$-measurable), we definitively extract it directly from the conditional expectation:
    \begin{equation}
        \frac{S_{t-1}}{(1+r)^t} \mathbb{E}^\mathbb{Q} \left[ 1+R_t \mid \mathcal{F}_{t-1} \right] = \frac{S_{t-1}}{(1+r)^{t-1}}.
    \end{equation}
    Safely dividing $S_{t-1}$ from both symmetrical sides (given $S_{t-1} > 0$ strictly) and clearing the discounting term algebraically:
    \begin{equation}
        \mathbb{E}^\mathbb{Q} \left[ 1+R_t \mid \mathcal{F}_{t-1} \right] = 1+r \implies \mathbb{E}^\mathbb{Q} \left[ R_t \mid \mathcal{F}_{t-1} \right] = r.
    \end{equation}
    To extract profound intuition, expand this explicit expectation forcefully across the two finite states of $R_t \in \{u, d\}$. Defining conditionally $q_t(\omega) = \mathbb{Q}[R_t=u \mid \mathcal{F}_{t-1}](\omega)$. We reduce to:
    \begin{equation}
        q_t u + (1-q_t) d = r.
    \end{equation}
    Resolving strictly to locate $q_t$ evaluates rapidly:
    \begin{equation}
        q_t(u - d) = r - d \implies q_t = \frac{r-d}{u-d}.
    \end{equation}
    We powerfully note that the analytically derived probability $q_t$ manifests purely as a deterministic constant! It completely avoids inheriting any variance from historical pathing dependencies via $\mathcal{F}_{t-1}$ and remains static across all structural times $t$. We categorically conclude that under $\mathbb{Q}$, isolated one-period returns actively operate wholly independent of preceding history, proving their explicit i.i.d. nature securely under $\mathbb{Q}$. Conditioned purely upon securely maintaining $d < r < u$, the mathematical parameter guarantees boundaries strictly holding as $q \in (0,1)$.
    
    Given $\mathbb{Q}$ possesses solitary unique status (which mathematically locks down the underlying completeness theorem natively, definitively empowering arbitrary payoff synthesis), we subsequently resolve that sequential transitions for strictly recursive variables heavily matching $S_t = S_{t-1}(1+R_t)$ function with flawless Markov property under $\mathbb{Q}$.
\end{proof}

\chapter{Pricing of Derivatives}

\section{The General Pricing Formula}

In a financial market, a derivative (or contingent claim) is a contract whose payoff depends on the behavior of the underlying primary assets. Mathematically, a European derivative expiring at maturity $T$ is defined as an $\mathcal{F}_T$-measurable random variable $H$. The $\mathcal{F}_T$-measurability ensures that the payoff $H$ is completely determined by the information available in the market up to time $T$.

In the context of our binomial model with one risky asset, the payoff of any derivative can be expressed as a deterministic function $h$ of the entire price path:
\begin{equation}
    H = h(S_0, S_1, \dots, S_T).
\end{equation}

According to the fundamental precepts of risk-neutral valuation, in an arbitrage-free and complete market, the unique no-arbitrage price of the derivative at any time $t \in \{0, \dots, T\}$, denoted by $V_t$, is given by the conditional expectation of the discounted payoff under the unique equivalent martingale measure $\mathbb{Q}$:
\begin{equation}
    V_t = \mathbb{E}^\mathbb{Q} \left[ \frac{H}{(1+r)^{T-t}} \;\middle|\; \mathcal{F}_t \right].
\end{equation}

Since the filtration $\mathcal{F}_t$ is generated by the stock prices up to time $t$, we intuitively expect the price $V_t$ to depend only on the observed path $(S_0, \dots, S_t)$. The following theorem formalizes this using the properties of the binomial model.

\begin{theorem}[Pricing Function for Path-Dependent Derivatives]
    Under the binomial model, the value of the replicating portfolio for a derivative with payoff $H = h(S_0, S_1, \dots, S_T)$ at time $t$ is given by:
    \begin{equation}
        V_t = v_t(S_0, S_1, \dots, S_t), \quad t = 0, \dots, T,
    \end{equation}
    where the deterministic pricing function $v_t : \mathbb{R}^{t+1} \to \mathbb{R}$ is defined explicitly by replacing the unknown future returns with their random variables under $\mathbb{Q}$:
    \begin{align}
        v_t(x_0, \dots, x_t) &:= \frac{1}{(1+r)^{T-t}} \times \nonumber \\
        &\quad \mathbb{E}^\mathbb{Q} \left[ h\left(x_0, \dots, x_t, x_t(1+R_{t+1}), \dots, x_t \prod_{k=t+1}^T (1+R_k)\right) \right]. \label{eq:general_vt}
    \end{align}
\end{theorem}

\begin{proof}
    By the risk-neutral pricing formula, $V_t = \frac{1}{(1+r)^{T-t}} \mathbb{E}^\mathbb{Q}[h(S_0, \dots, S_T) \mid \mathcal{F}_t]$. 
    We can rewrite the future prices $S_{t+1}, \dots, S_T$ in terms of the time-$t$ price $S_t$ and the future one-period returns $R_{t+1}, \dots, R_T$:
    \begin{equation}
        S_k = S_t \prod_{i=t+1}^k (1+R_i), \quad \text{for } k > t.
    \end{equation}
    Therefore, the payoff $H$ becomes:
    \begin{equation}
        H = h\left(S_0, \dots, S_t, S_t(1+R_{t+1}), \dots, S_t \prod_{k=t+1}^T (1+R_k)\right).
    \end{equation}
    Observe that $(S_0, \dots, S_t)$ are structurally $\mathcal{F}_t$-measurable (they are known constants given the information at time $t$), while the future returns $(R_{t+1}, \dots, R_T)$ are independent of $\mathcal{F}_t$ under the measure $\mathbb{Q}$ (since returns are i.i.d. under $\mathbb{Q}$). 
    
    By the Independence Lemma for conditional expectations, if we have a function $g(X, Y)$ where $X$ is $\mathcal{F}_t$-measurable and $Y$ is independent of $\mathcal{F}_t$, then $\mathbb{E}[g(X,Y) \mid \mathcal{F}_t] = \phi(X)$ where $\phi(x) = \mathbb{E}[g(x, Y)]$.
    Applying this powerful lemma, we substitute dummy variables $(x_0, \dots, x_t)$ for the known path $(S_0, \dots, S_t)$ and take the unconditional expectation over the remaining independent random variables to recover the exact deterministic function $v_t$ provided in equation \eqref{eq:general_vt}.
\end{proof}

\section{Pricing Vanilla Options}

A European derivative is called path-independent or a "vanilla" option if its payoff depends solely on the terminal asset price $S_T$, rather than the entire historical trajectory. In this case, $H = h(S_T)$ for some function $h : \mathbb{R} \to \mathbb{R}$ (for instance, $h(x) = \max(x-K, 0)$ for a European Call).

By the Markov property of the binomial price process under $\mathbb{Q}$, the price of a vanilla option at time $t$ depends strictly upon the current price $S_t$, allowing us to dramatically simplify $v_t(S_0, \dots, S_t)$ to merely $v_t(S_t)$.

\begin{proposition}[Cox-Ross-Rubinstein Formula for Time 0 Pricing]
    Let $H = h(S_T)$ be a path-independent payoff. Its fundamental no-arbitrage price at inception $t=0$ is structurally quantified by:
    \begin{equation}
        v_0(S_0) = \frac{1}{(1+r)^T} \sum_{k=0}^T \binom{T}{k} q^k (1-q)^{T-k} h\left(S_0(1+u)^k(1+d)^{T-k}\right).
    \end{equation}
\end{proposition}

\begin{proof}
    At time $t=0$, the risk-neutral formula declares $v_0(S_0) = \frac{1}{(1+r)^T} \mathbb{E}^\mathbb{Q}[h(S_T)]$. Because $S_T = S_0 \prod_{i=1}^T (1+R_i)$, the terminal price depends entirely upon the total number of "up" movements that unfold dynamically over the $T$ periods. 
    
    Let $X$ denote the absolute number of "up" movements (where $R_i = u$). Since the returns $R_i$ are independent under $\mathbb{Q}$ and identically yield $u$ with strict probability $q$, the variable $X$ intrinsically follows a Binomial distribution: $X \sim \mathcal{B}(T, q)$.
    If $X = k$, then there are exactly $T-k$ "down" movements, driving the terminal price flawlessly to $S_0(1+u)^k(1+d)^{T-k}$.
    Taking the rigorous mathematical expectation over the explicit probability mass function of the binomial distribution directly yields the weighted summation formula above.
\end{proof}

This explicit formula is highly practically efficient, requiring only $\mathcal{O}(T)$ operations. However, to track the price path dynamically through the tree or to price American options, backward induction stands as the preferred algorithmic strategy.

\begin{proposition}[Backward Induction Algorithm]
    For a generalized derivative, the deterministic pricing function $v_t$ thoroughly satisfies the recursive one-step relation functionally bridging time $t$ to $t+1$:
    \begin{align} \label{eq:backward_induction}
        v_t(S_0, \dots, S_t) &= \frac{1}{1+r} \Big[ q v_{t+1}(S_0, \dots, S_t, S_t(1+u)) \nonumber \\
        &\quad + (1-q) v_{t+1}(S_0, \dots, S_t, S_t(1+d)) \Big]
    \end{align}
    subject to the strict terminal boundary condition at maturity $T$:
    \begin{equation}
        v_T(S_0, \dots, S_T) = h(S_0, \dots, S_T).
    \end{equation}
\end{proposition}

\begin{proof}
    This powerful relation flows seamlessly from the innate tower property of conditional expectations. Specifically:
    \begin{equation}
        V_t = \mathbb{E}^\mathbb{Q} \left[ \frac{V_{t+1}}{1+r} \;\middle|\; \mathcal{F}_t \right].
    \end{equation}
    Since $V_{t+1} = v_{t+1}(S_0, \dots, S_t, S_{t+1})$ and $S_{t+1} = S_t(1+R_{t+1})$, we condition entirely upon the $\mathcal{F}_t$-measurable sequence $(S_0, \dots, S_t)$. The single element rigidly remaining uncertain conditionally represents $R_{t+1}$. Under the structural paradigm dictated by $\mathbb{Q}$, $R_{t+1}$ equals $u$ with explicit probability $q$ mapping the future path toward state $(S_0, \dots, S_t, S_t(1+u))$, or rigidly falls to $d$ with probability $(1-q)$ navigating heavily to $(S_0, \dots, S_t, S_t(1+d))$. Taking the expectation over this strictly isolated one-period jump identically produces the recursive algorithm in \eqref{eq:backward_induction}.
\end{proof}

\chapter{Replication of Derivatives}

\section{The Concept of Replication}

A core tenet of modern financial theory asserts that in a complete market, any derivative can be perfectly hedged. This means developing a dynamic trading strategy in the primary assets (the stock and the bank account) that miraculously synthesizes the exact payoff of the derivative under all possible future scenarios. Such a strategy is called a \textit{replicating portfolio}.

In mathematical terms, a trading strategy is defined by a pair of predictable processes $(\xi_t, \xi_t^0)_{t=1}^T$. 
\begin{itemize}
    \item $\xi_t$ represents the number of shares of the risky asset held during the period $(t-1, t]$.
    \item $\xi_t^0$ represents the number of units of the riskless asset (bank account) held during the period $(t-1, t]$.
\end{itemize}
Predictability requires that the portfolio components $\xi_t$ and $\xi_t^0$ must be rigorously determined based strictly upon the information available at the beginning of the period, namely $\mathcal{F}_{t-1}$. 

The value of this portfolio at time $t$ is strictly given by:
\begin{equation}
    V_t = \xi_t S_t + \xi_t^0 S_t^0.
\end{equation}
A derivative with an $\mathcal{F}_T$-measurable payoff $H$ is considered replicable if there exists a self-financing strategy such that its terminal portfolio value exactly logically matches the payoff almost surely:
\begin{equation}
    V_T = \xi_T S_T + \xi_T^0 S_T^0 = H \quad \mathbb{P}\text{-a.s.}
\end{equation}

\section{Dynamic Delta Hedging in the Binomial Model}

Because our discrete binomial model branches into exactly two unique states at every structural node, we possess exactly two independent financial instruments (the risky stock and the riskless bond) to balance the equation. This yields a mathematically determined $2 \times 2$ linear system of equations at every node guaranteeing perfect replication.

\begin{proposition}[Replication Strategy for Path-Dependent Payoffs]
    The perfectly self-financing replication strategy for a path-dependent European payoff $H = h(S_0, \dots, S_T)$ is determined at every node $t \in \{1, \dots, T\}$ by:
    \begin{equation}
        \xi_t = \frac{v_t(S_0, \dots, S_{t-1}, S_{t-1}(1+u)) - v_t(S_0, \dots, S_{t-1}, S_{t-1}(1+d))}{S_{t-1}(u-d)} \label{eq:delta_hedging}
    \end{equation}
    and
    \begin{equation}
        \xi_t^0 = \frac{v_t(S_0, \dots, S_{t-1}, S_{t}) - \xi_t S_t}{S_t^0}. \label{eq:bank_account}
    \end{equation}
\end{proposition}

\begin{proof}
    Consider the portfolio transitioning functionally from time $t-1$ across the single period to time $t$. Given the entire explicit path history up to $t-1$, strictly denoted conceptually as $\vec{S}_{t-1} = (S_0, \dots, S_{t-1})$, the underlying asset price rigidly moves to either $S_t^{up} = S_{t-1}(1+u)$ or $S_t^{down} = S_{t-1}(1+d)$. 
    
    Correspondingly, the target derivative theoretical price seamlessly transitions to either $V_t^{up} = v_t(\vec{S}_{t-1}, S_t^{up})$ or $V_t^{down} = v_t(\vec{S}_{t-1}, S_t^{down})$.
    
    To violently secure complete replication, our holding strategy $(\xi_t, \xi_t^0)$ decided cleanly at time $t-1$ absolutely must financially output the target value in both disparate eventualities. We construct the simultaneous system:
    \begin{align}
        \xi_t S_{t-1}(1+u) + \xi_t^0 S_t^0 &= V_t^{up} \label{eq:sys1} \\
        \xi_t S_{t-1}(1+d) + \xi_t^0 S_t^0 &= V_t^{down} \label{eq:sys2}
    \end{align}
    Note that the riskless asset $S_t^0$ structurally holds identical deterministic value $(1+r)^t$ violently irrespective of the "up" or "down" shift in the stock. By aggressively algebraic subtraction of equation \eqref{eq:sys2} securely from \eqref{eq:sys1}, we strictly eliminate the bond holding entirely:
    \begin{equation}
        \xi_t S_{t-1}(1+u) - \xi_t S_{t-1}(1+d) = V_t^{up} - V_t^{down}.
    \end{equation}
    Factoring rigorously:
    \begin{equation}
        \xi_t S_{t-1} (1 + u - 1 - d) = \xi_t S_{t-1}(u-d) = V_t^{up} - V_t^{down}.
    \end{equation}
    Isolating mathematically for $\xi_t$ successfully recovers the canonical discrete delta-hedging ratio formally declared in \eqref{eq:delta_hedging}:
    \begin{equation}
        \xi_t = \frac{V_t^{up} - V_t^{down}}{S_{t-1}(u-d)}.
    \end{equation}
    
    Once the strictly necessary stock allocation $\xi_t$ is powerfully quantified, we algebraically resubstitute directly back into either generating equation (or neutrally into the generic relation $V_t = \xi_t S_t + \xi_t^0 S_t^0$) to conclusively govern the riskless asset position formally shown universally in \eqref{eq:bank_account}.
\end{proof}

The quantity $\xi_t$ fundamentally expresses the discrete sensitivity of the derivative's intrinsic price relative to the underlying stock price shifts. This acts as the rigorous discrete-time analog to the continuous "Delta": $\Delta_t = \frac{\partial V}{\partial S}$. By forcefully holding exactly $\xi_t$ localized units, the investor dynamically neuters any structural directional exposure to the risky variable.

\section{Simplification for Vanilla Options}

Just as path-independent payoffs yield exceptionally simplified functional pricing expressions, their required replication sequences equivalently collapse heavily in sheer dimensionality.

\begin{proposition}[Vanilla Hedging Compression]
    If the explicitly defined derivative payoff relies absolutely cleanly upon terminal value, $H = h(S_T)$, then the theoretical pricing dependencies permanently transition identically conforming to:
    \begin{equation}
        v_t(S_0, S_1, \dots, S_t) \equiv v_t(S_t).
    \end{equation}
    Consequently, the localized delta-hedging allocations seamlessly collapse down to strictly demanding:
    \begin{equation}
        \xi_t \equiv \xi_t(S_{t-1}) \quad \text{and} \quad \xi_t^0 \equiv \xi_t^0(S_{t-1}).
    \end{equation}
\end{proposition}

\begin{proof}
    Because structural returns $R_t$ naturally maintain their strict i.i.d status beneath the measure $\mathbb{Q}$, the intrinsic price process $S_t$ operates rigorously as a structurally complete Markov process. Consequently, conditionally analyzing any functionally prospective path-independent payoff $h(S_T)$ relies structurally purely upon realizing exactly the localized state variable mathematically isolated as $S_{t-1}$. Therefore, replacing heavily burdened path vectors inside equations \eqref{eq:delta_hedging} smoothly strips out dimensional excess, conclusively proving that hedging mechanics algorithmically demand strictly present-state awareness exclusively.
\end{proof}

\chapter{Structure of Complete Market}

\section{General Discrete-Time Markets}

Up until this point, we have meticulously analyzed a highly specialized discrete market containing exactly $d=1$ risky asset and $1$ riskless asset, constrained to exhibit strictly binomial price transitions. A natural theoretical question arises: Are there more generalized discrete-time frameworks that naturally yield complete markets? 

To rigorously investigate this structural property, consider a generalized $T$-period market modeled upon a probability space $(\Omega, \mathcal{F}, \mathbb{P})$. The market contains one riskless asset serving cleanly as the numeraire, and an arbitrary multidimensional vector of $d$ risky assets modeled collectively as a $d$-dimensional stochastic process $S = (S^1, \dots, S^d)$. 

The explicit transition dynamics of these multiple assets follow:
\begin{equation}
    \Delta S_t := S_t - S_{t-1} \in \mathbb{R}^d.
\end{equation}
The underlying multidimensional market is \textit{complete} if and only if every bounded $\mathcal{F}_T$-measurable random variable can be replicated via a self-financing strategy across all $d$ risky assets and the singular bank account unit. Driven by the Second Fundamental Theorem of Asset Pricing, this holds if and exclusively if the equivalent martingale measure $\mathbb{Q}$ acts uniquely.

\section{The \texorpdfstring{$d+1$}{d+1} Branching Theorem}

The geometric constraints governing market completeness lock branching probability distributions toward low dimensional geometries matching the underlying dimension of liquid financial instruments.

\begin{proposition}[Geometric Branching Structure of Complete Markets]
    Assume that there exists a \textbf{unique} martingale measure (the market is arbitrage-free and complete). 
    
    Then, there exists an adapted process consisting of exactly $d+1$ vectors $\{\delta_{t}^i\}_{i=1}^{d+1}$ formally taking values in $\mathbb{R}^d$, tightly enforcing computationally that:
    \begin{equation} \label{eq:branching_condition}
        \sum_{i=1}^{d+1} \mathbb{P}\left(\Delta S_t = \delta_{t-1}^i \;\middle|\; \mathcal{F}_{t-1}\right) = 1 \quad \mathbb{P}\text{-a.s.} \;\; \forall t \in \{1, \dots, T\}.
    \end{equation}
    Moreover, for all $t$, the family of state vectors $\{\delta_t^1, \dots, \delta_t^{d+1}\}$ is strictly \textbf{affinely independent} $\mathbb{P}$-a.s. Mathematically, this enforces that the derivative vectors:
    \begin{equation}
        \{ \delta_t^2 - \delta_t^1, \dots, \delta_t^{d+1} - \delta_t^1 \}
    \end{equation}
    are \textbf{linearly independent} vectors strongly spanning $\mathbb{R}^d$ almost surely.
\end{proposition}

\begin{proof}
    The geometric mappings dominating multi-asset pricing heavily rely upon balancing the equations dynamically. At any node, if the path branches into $K$ disjoint possibilities, we generate $K$ target portfolio values to match. Our strategy permits choosing explicit allocations spanning exactly $d$ risky assets supplemented alongside $1$ riskless bank account. 
    Therefore, perfect replication yields a system of $K$ equations with exactly $d+1$ structural unknowns. Resolving this unique geometric parity demands $K \leq d+1$. The linear independence rigidly ensures full spatial span to prevent redundant paths, conclusively enforcing uniquely solvable algorithmic replication.
\end{proof}

\section{The Singularity of the Binomial Model}

This abstract spatial theorem immediately generates powerful analytical consequences regarding discrete-time modeling constraints.

Substituting identical values for $d=1$ smoothly generates exactly $1+1=2$ explicitly valid branching paths. Thus:
\begin{enumerate}
    \item A non-trivial model prevents pure market stagnation, generating bounded price shifts: $\mathbb{P}[S_t = S_{t-1}] < 1$.
    \item The functionally necessary matching of exactly $2$ independent branches fiercely limits complete 1-dimensional modeling flawlessly to the Cox-Ross-Rubinstein \textbf{binomial model}.
\end{enumerate}

\begin{corollary}[Incompleteness of Complex Trees]
    Consequently, excluding isolated tree arrays precisely configured to span $d+1$ branches for exactly $d$ independent liquid assets, generalized discrete time structural models naturally and decisively exhibit pure \textbf{incompleteness}.
\end{corollary}

In explicit financial vocabulary, heavily attempting to utilize trinomial (3-branch) trees to map completely hedgeable options operating upon merely a single ($d=1$) underlying asset predictably fails. The excess unknown states relative to the explicit trading tools natively embeds unhedgable risk structurally operating fundamentally beyond replication entirely.


\end{document}
